\documentclass[10pt,a4paper]{amsart}
\usepackage[margin=19mm]{geometry}
\usepackage{amsmath,amssymb,mathtools}
\usepackage[scr=rsfs,cal=euler]{mathalfa}
\usepackage{xfrac}
\usepackage[unicode,hidelinks,bookmarks=false]{hyperref}
\newcommand{\ie}{i.e.\ }
\newcommand{\cf}{cf.\ }
\newcommand{\resp}{respectively\ }
\newcommand{\Subsection}{\subsection}
\providecommand{\nobreakdash}{\nobreak\hbox{-}}
\setlength{\emergencystretch}{2em}
\multlinegap0pt
\usepackage{amssymb}
\usepackage{bm}
\newtheorem{proposition}{Proposition}
\newtheorem{theorem}{Theorem}
\title{Locally Repairable Codes with Availability via Elliptic Function Fields}
\author{Junjie Huang, Chang-An Zhao}
\date{Source: arXiv:2605.06182v1 (CC BY 4.0)}
\begin{document}
\maketitle

\noindent\emph{Source and licence.} Locally Repairable Codes with Availability via Elliptic Function Fields. Version arXiv:2605.06182v1 (CC BY 4.0), distributed by arXiv under the Creative Commons Attribution 4.0 International licence, http://creativecommons.org/licenses/by/4.0/, as recorded in the arXiv licence metadata for this exact version and shown on the abstract page https://arxiv.org/abs/2605.06182v1. Full licence notice: \texttt{licenses/arxiv-2605.06182-CC-BY-4.0.txt}.\par\smallskip

\noindent\textit{Editorial scope.} Counted unit: the article's theorem theorem (source0.tex lines 909--919). The article's complete original proof of it is delivered with it (source0.tex lines 920--930, one \texttt{proof} environment of 11 lines). No mathematics is written, repaired or re-derived: the statement and the proof are the article's own lines, in the article's own notation, and no retained line is substituted or re-flowed.  Retained uncounted as the source-local context the proof needs: lines 412--418; lines 782--791; lines 863--883.  Nothing was omitted from any retained span, so the package carries no omission record.  No retained line contains a citation, so the wrapper carries no bibliography and no bibliography entry.  No retained line contains a figure environment, an image-inclusion command or drawing code, so no image is delivered and the package declares no asset dependency.  Editorial formatting only, in addition: an AMS \texttt{amsart} preamble that restates the article's own declarations and macros used by the retained lines, namely the environments \texttt{proposition}, \texttt{theorem} on separate counters, together with the standard packages those lines need.  The retained environments print in this document as Proposition 1, Theorem 1 in order of appearance, whereas the article prints the same items with the numbering of its own full document.  The complete native source of the article as distributed in the arXiv e-print for this exact version is bundled unchanged as \texttt{sources/200003/source0.tex}.
\begin{proposition}\label{new_ellp_cuv}
    Let \( p \) be a  prime number.
    \begin{itemize}
        \item[(i)] Let \( \mathfrak{E}/\mathbb{F}_{p^2} \) be an elliptic curve over \( \mathbb{F}_{p^2} \) defined by the equation \( y^2 = x^3 + b \) for some \( b \in \mathbb{F}_{p}^* \) and \( E/\mathbb{F}_{p^2} \) be the elliptic function field of \( \mathfrak{E}/\mathbb{F}_{p^2} \). Then \(N(E) = {p^2} + 2p\) if and only if \(p = 3u^2 + 3u + 1\) with \(u\in\mathbb{Z}\). In this case, we have \(|\operatorname{Aut}(E, \mathcal{O})| = 6\). 
        \item[(ii)] Let \( \mathfrak{E}/\mathbb{F}_{p^2} \) be an elliptic curve over \( \mathbb{F}_{p^2} \) defined by the equation \( y^2 = x^3 + ax \) for some \( a \in \mathbb{F}_{p}^* \) and \( E/\mathbb{F}_{p^2} \) be the elliptic function field of \( \mathfrak{E}/\mathbb{F}_{p^2} \). Then \(N(E) = {p^2} + 2p - 3\) if and only if \( p \equiv 1 \pmod{4} \) and \(p = v^2 + 1\) with \(v\in\mathbb{Z}\). In this case, we have \(|\operatorname{Aut}(E, \mathcal{O})| = 4\). 
    \end{itemize}
\end{proposition}

\begin{proposition}\label{cons_lrc_two}
    Let \( E/\mathbb{F}_q \) be an elliptic function field. Let \(H = \{P_1,P_2,\cdots,P_{|H|} = \mathcal{O}\}\) be a subgroup of \( \mathbb{P}^1_E \) with \( |H|<q \). Let \( T_H = \{\tau_P:P\in H\} \) be a subgroup of \( T_E \) and let \( A_1,A_2 \) be two nontrivial subgroups of \( \operatorname{Aut}(E, \mathcal{O}) \) such that \(|A_1\cap A_2| = 1\) and \( \mathcal{G} = T_H A_1 A_2 \) is a subgroup of \( \operatorname{Aut}(E/\mathbb{F}_q) \). Suppose that \( Q_1, \cdots, Q_m \) are rational places of \( E^{\mathcal{G}} \) that split completely in \( E \) and the rational place $P_{i,j}$ of $E$ lying over $Q_i$ satisfies that \(P_{i,j}\notin E[|H|]\) for any $1\le i\le m$ and \(1\le j\le |H||A_1||A_2|\). Put \( r_1 = |H||A_1| - 1 \), \( r_2 = |H|(|A_2| - 1)\) and \( n = m|H||A_1||A_2| \). Then for any integer \( d_0 \) with \( 1 \leq d_0 < n \), there exists a \( q \)-ary \([n, k, d; (r_1, r_2)]_q\) locally repairable code with \(d \ge n - L\ge d_0\) and  
    \[
        k \geq \bigg\lfloor\frac{n-d_0}{r_1 + 1}\bigg\rfloor r_1 + \bigg\lfloor\frac{n-d_0}{r_2 + |H|}\bigg\rfloor r_2  + 2 - L,
    \]
    where
    \[
        L =  \max\left\{\left\lfloor\frac{n-d_0}{r_1 + 1}\right\rfloor (r_1+1),\left\lfloor\frac{n-d_0}{r_2 + |H|}\right\rfloor (r_2 + |H|)\right\}.  
    \]
\end{proposition}

\begin{theorem}\label{lrc_two_max_restate}
    Let \( q = p^a \) for any prime \( p \) and any even integer \( a > 0 \). Let \(\sqrt{q}+1 = \prod_\ell \ell^{h_\ell}\)  be the prime factorization of \( \sqrt{q} + 1 \). For any integer \( d_0 \) with \( 1 \leq d_0 < n \) and any positive divisor \( h \) of \( \sqrt{q} + 1 \) with
    \begin{equation}\label{lrc_two_repair_condition}
        h^2|A_1||A_2|>\prod_\ell \ell^{2\min\{2\nu_\ell(h),h_\ell\}}-h^2,
    \end{equation}
    there exists a \( q \)-ary \([n = mh^2|A_1||A_2|, k, d;(r_1,r_2)]_q\) locally repairable code where \( r_1 = h^2|A_1| - 1 \), \(r_2 = h^2(|A_2|-1)\), \(d \ge n - L \ge d_0\), 
    \[
        k \geq \bigg\lfloor\frac{n-d_0}{r_1 + 1}\bigg\rfloor r_1 + \bigg\lfloor\frac{n-d_0}{r_2 + h^2}\bigg\rfloor r_2  + 2 - L,
    \]
    and
    \[
        L = \max\left\{\left\lfloor\frac{n-d_0}{r_1 + 1}\right\rfloor (r_1+1),\left\lfloor\frac{n-d_0}{r_2 + h^2}\right\rfloor (r_2 + h^2)\right\}, 
    \]
    for any integer \( m \) satisfying \(1 \leq m \leq \left\lceil \frac{q+2\sqrt{q}+1-2h^2}{h^2|A_1||A_2|} \right\rceil - 1 \), provided that \(|A_1|\), \(|A_2|\) and \( p \) satisfy one of the following cases:
     
    \begin{itemize}
        \item[(i)] \(|A_1| = 2,4,8\) and \(|A_2| = 3\) for \( p = 2 \);  
        \item[(ii)] \(|A_1| = 2,4\) and \(|A_2| = 3\) for \( p = 3 \);  
        \item[(iii)] \(|A_1| = 2\) and \(|A_2| = 3\) for \( p \equiv 2 \pmod{3} \) and \( p \neq 2 \).
    \end{itemize}
\end{theorem}

\begin{theorem}
    Let \( q \) be a prime power. For any integer \( d_0 \) with \( 1 \leq d_0 < n \) and any positive divisor \( h \) with \(h\mid\mid q^2+2q\), there exists a \( q^2 \)-ary \([n = mh|A_1||A_2|, k, d;(r_1,r_2)]_{q^2}\) locally repairable code where \( r_1 = h|A_1| - 1 \), \(r_2 = h(|A_2|-1)\), \(d \ge n - L \ge d_0\), 
    \[
        k \geq \bigg\lfloor\frac{n-d_0}{r_1 + 1}\bigg\rfloor r_1 + \bigg\lfloor\frac{n-d_0}{r_2 + h}\bigg\rfloor r_2  + 2 - L,
    \]
    and
    \[
        L = \max\left\{\left\lfloor\frac{n-d_0}{r_1 + 1}\right\rfloor (r_1+1),\left\lfloor\frac{n-d_0}{r_2 + h}\right\rfloor (r_2 + h)\right\} , 
    \]
     for any integer \( m \) satisfying \(1 \leq m \leq \left\lceil \frac{q^2+2q-2h}{h|A_1||A_2|} \right\rceil - 1 \), provided that \(|A_1|\), \(|A_2|\) and \( q \) satisfy that \(|A_1| = 2\) and \(|A_2| = 3\) for \(q = 3u^2 + 3u + 1\) with \(u\in\mathbb{Z}\). 
\end{theorem}

\begin{proof}
    Apart from using the elliptic curve in Proposition \ref{new_ellp_cuv}(i) and the corresponding elliptic function field, the proof is analogous to that of Theorem \ref{lrc_two_max_restate}. 
    
    Let \(H = E[h]\). Then, by Proposition \ref{new_ellp_cuv}(i), we can derive \(|H| = h\) and the existence of subgroups \(T_H\), \(A_1\) and \(A_2\) that satisfy the above assumptions together with the conditions in Proposition \ref{cons_lrc_two}. Therefore, we only need to prove that for any rational place \(Q\) of \( E^{T_HA_1A_2} \) that split completely in \( E \), the rational place $P$ of $E$ lying over $Q$ satisfies that \(P\notin E[|H|]\); then, together with Proposition \ref{cons_lrc_two}, we can obtain our conclusion. 

    Suppose, on the contrary, that \(P\in E[|H|]\). Then we can also obtain 
    \[
        h|A_1||A_2|\le |E[|H|]|-|E[h]|.
    \]
    However, it can be readily verified that \(|E[|H|]|=|E[h]|=h\). So we have \(h|A_1||A_2|\le 0 \) which is a contradiction. Hence, we have \(P\notin E[|H|]\). The proof is completed. 
\end{proof}

\end{document}
